\clearpage
\section*{40. 非可換代数}

\begin{center}
\emph{Math. Zs. 37 (1933), 514--541}
\end{center}

可換代数の主要定理は、よく知られているように、ガロア理論に含まれており、その前段には添加体と分解体の理論が置かれる。すなわち、与えられた多項式が一次因子を一つ分離するため、または完全に一次因子へ分解するために十分な体の理論である。

ここではそれに対応する代数の部分を、非可換の場合、とくに超複素の場合に展開する。原理的には非可換的方法、すなわち非可換体（斜体）における表現を用いる。最後に、上に述べた可換代数の定理が、可換体における表現によって完全に平行な仕方で証明されることを示す。

用いられる表現論は、短い自己同型論（§1）を前置したうえで、表現加群に基づく理論をさらに発展させたものである。とくに、直接表現のほかに相反表現を考える。両者は互いに還元できるが、ここでは相反表現は相反表現加群によって生成される。この利点は、その相反表現加群が、したがって双加群が、拡大環上の片側加群としても見られる点にある（§2）。こうして斜体における表現は拡大環におけるイデアル論に還元される。既約表現類は拡大環の既約イデアル類に対応し、これは超複素系そのものを可換な係数体上で表現したときに成り立つ事実と正確に同じである（§§3, 4）。

この立場から、非可換体上の行列環についての構造定理が得られる。すなわち、任意の部分環はこの体による表現を、同値にいえば相反同型な体による相反表現を与える、という観察（§5）による。この相反同型な体は、可換の場合における最小添加体および分解体の第一の類似物である。これはすべての一次の相反表現を媒介し、また中心上の階数が有限の場合には階数 \(1\) の直和因子への完全分解も媒介する。これにより斜体のガロア理論（§6）が得られ、さらに、相反同型な除法代数、すなわち中心上有限階数の体が、R. Brauer の代数類群において逆類を生成するという事実（§7）も表される。

より一般に単純系に有効なガロア理論の第二の証明は、可換化する部分環に関する定理のほとんど直接の帰結であり、この定理自体も前の自己同型に関する考察とほぼ直接につながっている。ここまでは議論は純粋に非可換である。しかし可換な分解体の問題も、§5 で準備した観察を用い、分解体を除法代数を通じて表現することにより非可換的に扱われる（§7）。最後の部分では可換の場合への移行（§8）と任意の系に対する分解体の理論（§9）を与える。

R. Brauer は分解体の理論を、可換の場合のこの表現論と非有理的な因子系の上に築いた。この可換的基礎のために、彼、そして後に Albert は、中心が完全体であることを仮定しなければならなかったが、この制限は非可換的基礎では不要である。同じ制限のもとで、やはり可換の場合の表現論を用いて、R. Brauer と K. Shoda は理論をさらに発展させた。

最後に、斜体については、1928 年夏の私の講義に従った van der Waerden の短い叙述があることを述べておく。van der Waerden の叙述はその講義に対していくつかの簡略化を導入した。その一部は二度目の講義で採用し、一部はここで初めて採用した。可換化する部分環に関する定理（§5）と、それに基づく非可換ガロア理論の第二の証明（§6, 1 および 2）は、後になって加えられたものである。第一の講義におけるある議論、すなわち差分イデアルの交わりを作る方法は、より複雑ではあるが、無限階数の系および整な系に移行できるという利点を持つ。可換な整系については、これがイデアル微分とディフェレント（Different）との関係へ導く。

分解体の理論の主要な応用は、交叉積とその因子系の理論にあり、それはさらに数論への応用の基礎となる。ここではこれ以上追わない。

\subsection*{§1. 自己同型、加群、双加群について}

表現加群に基づく表現論は、作用素を持つ、または持たないアーベル群の自己同型環の理論に依拠する。そこで暗に用いられる議論、すなわち写像と計算法則との関係を、反復を避けるためにいくつかの単純な命題として定式化しておく。

\paragraph{1. 乗法的写像。結合法則。}
まず \(\frG\) を作用素を持たない群とし、\(\frA\) をその絶対自己同型領域、すなわち \(\frG\) から自身へのすべての準同型の系とする。\(\frA\) は乗法に関して閉じている。なぜなら積 \(\sigma\tau\) は
\[
        g(\sigma\tau)=(g\sigma)\tau,
        \qquad g\in\frG,\quad \sigma,\tau\in\frA,                       \tag{1}
\]
によって定義され、明らかに結合法則を満たすからである。

記号 \(O,H,\ldots\) からなる集合 \(\frB\) が与えられ、積 \(gO,gH,\ldots\) が \(\frG\) の元として定義され、かつ \(\frG\) から自身への準同型を生成するとき、群 \(\frG\) は作用素を持つ群となる：
\[
        (gh)O=gO\cdot hO.
\]
作用素領域 \(\frB\) が乗法的に閉じているなら、自己同型領域への誘導写像が乗法的準同型であることは、結合関係
\[
        g(OH)=(gO)H,
        \qquad g\in\frG,
        \ O,H\in\frB                                                   \tag{1a}
\]
が成り立つことと同値である。対応して、左作用素については相反準同型が得られる。

\paragraph{2. 作用素準同型写像。}
\(\frG\) が作用素を持つ群であるとき、作用素準同型性は
\[
        (gO)\sigma=(g\sigma)O,                                          \tag{2}
\]
によって、左作用素については
\[
        (Og)\sigma=O(g\sigma)                                           \tag{2*}
\]
によって定義される。二つの作用素領域 \(\frB\) と \(\frC\) について、\(\frC\) の元が \(\frB\)-自己同型を生成し、同時に \(\frB\) の元が \(\frC\)-自己同型を生成するのは、それらの領域が \(\frG\) と可換的に結びついているとき、すなわち走る結合法則
\[
        (gO)H=(gH)O,                                                    \tag{2a}
\]
または
\[
        (OH)g=O(Hg)                                                     \tag{2a*}
\]
が満たされるとき、かつそのときに限る。

\paragraph{3. 環上の加群と双加群。}
環 \(\frR\) 上の右加群 \(\frM\) とは、\(\frR\) の元を右作用素として持ち、結合関係と分配関係
\[
        (g+h)\rho=g\rho+h\rho,
        \qquad
        g(\rho+\sigma)=g\rho+g\sigma                                  \tag{3a}
\]
を満たす加法的アーベル群である。左加群についても対応する定義が成り立つ。

双加群には二種類ある。二つの環 \(\frR\) と \(\frS\) 上の右加群は、個別の加群法則に加えて
\[
        (m\rho)\sigma=(m\sigma)\rho
\]
を満たすとき双加群と呼ばれる。\(\frR\)-左、\(\frS\)-右加群は、走る結合法則
\[
        \rho(m\sigma)=(\rho m)\sigma
\]
を加えるとき双加群である。

加法的に書かれたアーベル群 \(\frM\) の作用素領域 \(\frR\) が環であるとき、\(\frR\) から絶対自己同型環への写像が環準同型であることは、\(\frM\) が右 \(\frR\)-加群であることとちょうど同値である。左加群の場合には相反的な環準同型になる。\(\frM\) が \(\frR,\frS\) 上の右加群であるとき、\(\frM\) が双加群であることは、\(\frR\) が \(\frS\)-自己同型環の部分環へ準同型に写り、\(\frS\) が \(\frR\)-自己同型環の部分環へ準同型に写ることとちょうど同値である。\(\frR\)-左、\(\frS\)-右双加群では、\(\frR\) の写像は相反であり、\(\frS\) の写像は直接である。

とくに次が成り立つ。

1) \(\frR\) が単位元を持つ環なら、\(\frR\) は左 \(\frR\)-加群としての自己同型環に直接同型であり、右 \(\frR\)-加群としての自己同型環に相反同型である。

2) \(\frR\) が一般には非可換な体 \(A\) 上の全行列環で、
\[
        \frR=\sum c_{ik}A=\sum A c_{ik},
\]
ならば、\(A\) は単純右イデアルの自己同型体に相反同型であり、単純左イデアルの自己同型体に直接同型である。

\paragraph{4. 右双加群から積環上の片側加群への移行。}
右双加群の意義は本質的にこの移行に基づく。

\paragraph{定理。}
\(\frM\) を、互いに元ごとに可換な二つの部分環 \(\frR\) と \(\frS\) を含む環 \(\frT\) 上の右加群とする。このとき \(\frM\) は \(\frR\) と \(\frS\) 上の双加群とも見なせる。逆に、\(\frR\) と \(\frS\) 上の双加群が与えられ、\(\frR\) と \(\frS\) が元ごとに可換となる積環 \(\frT\) が存在するなら、\(\frM\) は \(\frT\) 上の右加群とも見なせ、その \(\frT\)-作用は与えられた \(\frR\)-作用と \(\frS\)-作用を延長する。

絶対自己同型環の中の積環は、
\[
        \sum_i \rho_i\sigma_i+\rho+\sigma
\]
の形のすべての元からなる（\(\frR,\frS\) が単位元を持つなら、余分な項 \(\rho,\sigma\) は消える）。準同型
\[
        \frR\to\overline{\frR},\qquad
        \frS\to\overline{\frS}
\]
は
\[
        \sum_i \rho_i\sigma_i+\rho+\sigma
        \longmapsto
        \sum_i \overline\rho_i\,\overline\sigma_i+\overline\rho+\overline\sigma
\]
によって \(\frT\to\overline{\frT}\) の準同型へ延長される。したがって
\[
        m\Bigl(\sum_i \rho_i\sigma_i+\rho+\sigma\Bigr)
        =\sum_i (m\rho_i)\sigma_i+m\rho+m\sigma
\]
はよく定義され、与えられた \(\frR,\frS\)-加群構造を含む \(\frT\)-加群を定める。

\subsection*{§2. 相反表現と直接表現}

\paragraph{1. 線形形式の加群。}
\(\frS\) を単位元を持つ環とする。右 \(\frS\)-加群 \(\frM\) が \(\frS\) における線形形式の加群と呼ばれるのは、\(\frM\) が \(n\) 個の一項 \(\frS\)-加群の直和
\[
        \frM=m_1\frS+\cdots+m_n\frS
\]
であり、各 \(m_i\frS\) が \(\frS\) と作用素同型、したがって \(m_i s=0\) なら常に \(s=0\) となるときである。

\paragraph{定理 1.}
\(\frM\) が \(\frS\) における線形形式の加群なら、\(\frM\) の \(\frS\)-自己同型環 \(\frU\) は、\(\frS\) 上の \(n\) 行全行列環 \(\overline{\frU}\) に、単なる準同型ではなく相反同型である。

\(\frM\) の任意の \(\frS\)-自己同型 \(\alpha\) は
\[
        m_i\longmapsto \overline m_i=m_i\alpha
\]
によって完全に記述される。したがって
\[
        \sum_i m_i s_i\longmapsto
        \sum_i \overline m_i s_i
        =\sum_i (m_i\alpha)s_i
        =\sum_i (m_i s_i)\alpha.
\]
さらに
\[
        (m_1\alpha,\ldots,m_n\alpha)=(m_1,\ldots,m_n)A
\]
と書けば、対応 \(\alpha\mapsto A\) は一対一であり、
\[
        \alpha+\beta\mapsto A+B,
        \qquad
        \alpha\beta\mapsto BA.
\]
よってこの同型は相反である。従って次を得る。

\paragraph{定理 1'.}
\(\frS\) における次数 \(n\) の線形形式の加群の自己同型環は、\(\frS\) 上の \(n\) 行全行列環によって、同型で忠実な相反表現を許す。

\paragraph{2. 表現と表現加群。}
\(\frS\) における線形形式の加群 \(\frM\) は、\(\frR\) と \(\frS\) 上の双加群であり、同時に \(\frR\) と \(\frS\) の双方について右加群であるとき、\(\frR\) の \(\frS\) における相反表現加群と呼ばれる。\(\frM\) が双加群であり、同時に左 \(\frR\)-加群かつ右 \(\frS\)-加群であるとき、直接表現加群と呼ばれる。

\paragraph{定理 2.}
任意の相反表現加群、また直接表現加群は、それぞれ \(\frR\) の \(\frS\) における同値な相反表現、また直接表現の一つの類を生成し、そのようなすべての表現はこの仕方で得られる。

相反表現加群に対して、\(\frR\) は \(\frM\) の \(\frS\)-自己同型環の部分環へ直接準同型に写る。そこで定理 1' が相反表現を与える。逆に、相反表現を \(\frS\)-自己同型環への相反同型と合成すると、\(\frR\) からの直接準同型が得られ、したがって表現加群が得られる。直接の場合は準同型性を反転して得られる。\(\frR\) の直接表現は、\(\frR\) の相反環の相反表現であり、その逆も同様である。

\paragraph{3. 相反表現加群から拡大環上の加群への移行。}
表現加群に対して移行定理は次の形をとる。

\(\frM\) が、互いに元ごとに可換な二つの部分環 \(\frR\) と \(\frS\) を含む環 \(\frT\) 上の右加群であり、かつ \(\frS\) 上の線形形式の加群であるなら、\(\frM\) は \(\frR\) の \(\frS\) における相反表現加群とも見なせる。逆に、相反表現加群は、積環 \(\frT\) が存在する限り \(\frT\)-加群と見なせる。

\(\frM\) が \(\frR\) の \(\frS\) における相反表現加群であり、かつ \(\frT\) を積環とする \(\frT\)-加群でもあるなら、\(\frM\) と加群 \(\frN\) との \(\frR,\frS\)-同型は \(\frT\)-同型と同値である。従って、同型な \(\frT\)-加群の各類は \(\frR\) の \(\frS\) における相反表現類に対応し、逆も同様である。

\paragraph{可換化する行列に関する定理。}
\(\frR\mapsto\frR^*\) を \(\frR\) の \(\frS\) における次数 \(n\) の相反表現とし、\(\frB^*\) を \(\frS\) 上の \(n\) 行行列のうち \(\frR^*\) と元ごとに可換なすべてのものの環とする。このとき \(\frB^*\) は、それによって生成される表現加群の \(\frR,\frS\)-自己同型、すなわち同時に \(\frR\)-自己同型でもある \(\frS\)-自己同型の相反同型表現を与える。積環 \(\frT\) が存在すれば、それらは \(\frT\)-自己同型である。

\(\frR,\frS\)-自己同型は、同じ表現を生じる対応 \(m_i\mapsto m_i'\) にほかならない。元 \(m_i'\) は \(\frM\) の \(\frS\)-基底をなす必要はなく、したがって \(\frB^*\) の行列も可逆である必要はない。この拡張された意味では
\[
        (m_1',m_2',\ldots,m_n')a=(m_1',m_2',\ldots,m_n')A
\]
はなお正しい。ただし \(A\) はもはや \(a\) によって一意には決まらない。

対角行列 \(E\cdot s\) は \(\frS\) の直接同型な表現、すなわち \(\frS\) の恒等表現を与える。したがって \(\frR\) と \(\frS\) によって定義される \(\frT\) から \(\frS\) 上の行列への対応は、環 \(\frT\) の表現ではなく、\(\frR\) の相反表現を、\(\frS\) に関して作用素準同型なものへ延長したものである：
\[
        \sum_i r_i s_i\longmapsto \sum_i R_i s_i.
\]
とくに、\(\frR\) の相反表現は、\(\frR\) と \(\frS\) の交わり \([\frR,\frS]\) に関して作用素準同型であり、この交わりは \(\frR\) の中心に含まれる。

\subsection*{§3. 斜体に関する加群}

以下では、一般には非可換な斜体における表現だけを扱う。そのため、斜体上の線形形式の加群についての簡単な結果を集めておく。

\paragraph{1. 全加群の与えられた基底に関する部分加群の正規基底。}
\(A\) を斜体とし、
\[
        \frN=x_1A+\cdots+x_nA
\]
を階数 \(n\) の線形形式の加群とする。また
\[
        \frL=z_1A+\cdots+z_\ell A
\]
を階数 \(\ell\le n\) の部分加群とする。\(z_i\) が \(x_i\) に関する正規基底と呼ばれるのは、それらが
\[
        z_i=x_i-(x_{\ell+1}\alpha_{i,\ell+1}+\cdots+x_n\alpha_{i,n}),
        \qquad i=1,\ldots,\ell
\]
の形をしているときである。任意の加群 \(\frL\) は、\(x_i\) の適切な番号づけの後に、このような正規基底を持つ。実際、
\[
        \frN=\frL+x_{\ell+1}A+\cdots+x_nA
\]
であり、したがって
\[
        x_i\equiv x_{\ell+1}\alpha_{i,\ell+1}+\cdots+x_n\alpha_{i,n}
        \pmod{\frL},\qquad i=1,\ldots,\ell.
\]
従って対応する \(z_i\) は \(\frL\) に属し、\(\frL\) を尽くす。

\paragraph{2. 拡大加群。}
再び \(A\) を斜体とし、\(P\) を部分体とする。階数 \(n\) の \(A\)-加群 \(N\) が \(P\)-加群 \(M\) の拡大加群と呼ばれ、\(N=M_A\) と書かれるのは、\(M\) が \(N\) の同じ階数を持つ部分加群であるときである。もし
\[
        M=y_1P+\cdots+y_nP,
\]
なら
\[
        M_A=y_1A+\cdots+y_nA.
\]
逆に、任意の \(P\)-加群 \(M\) に対して、拡大加群は加群同型を除いて一意に存在する。

\paragraph{系 a).}
\(z_1,\ldots,z_s\) が \(M\) の中で \(P\) 上線形独立なら、それらは \(M_A\) の中で \(A\) 上も線形独立である。従って \(M\) の中の任意の \(P\)-加群
\[
        T=z_1P+\cdots+z_sP
\]
は拡大加群
\[
        T_A=z_1A+\cdots+z_sA\subset M_A
\]
を生成する。

\paragraph{系 b).}
\(M\) の中の任意の \(P\)-加群 \(T\) は収縮加群、すなわちその拡大加群と \(M\) との交わりである：
\[
        T=T_A\cap M .
\]

\paragraph{3. 不変加群に関する定理。}
\(N=M_A\) を \(P\)-加群 \(M\) の拡大加群とし、\(P\) を \(A\) の環自己同型群 \(\frG\) に対する完全な不変体とする。群 \(\frG\) は \(N\) 上の作用素領域として
\[
        Gm=m,
        \qquad m\in M,
\]
および
\[
        G\Bigl(\sum_i m_i\alpha_i\Bigr)=\sum_i m_i\,G(\alpha_i)
\]
により定義される。補題：\(M\) は \(\frG\) によって固定される \(N\) の元の全体である。

\paragraph{定理。}
\(M\) の部分加群 \(L\) の拡大加群 \(L_A\)、かつそれらだけが、\(\frG\) を作用素領域とする許容部分加群である。

一方向は明らかである。逆に、\(L\) を \(\frG\) に関して許容とし、\(z_1,\ldots,z_\ell\) を \(M\) の \(P\)-基底 \(x_1,\ldots,x_n\) に関する \(L\) の正規基底とする。任意の \(G\in\frG\) に対して
\[
        G(z_i)=x_i-(x_{\ell+1}G(\alpha_{i,\ell+1})+\cdots+x_nG(\alpha_{i,n}))
\]
は再び \(L\) に属する。\(x_1,\ldots,x_\ell\) の係数を比較すると \(G(z_i)=z_i\) を得る。従って補題により \(z_i\) は \(M\) に属し、\(L\) は \(M\) の中の \(P\)-加群
\[
        T=z_1P+\cdots+z_\ell P
\]
の拡大であり、\(T=L\cap M\) である。

\subsection*{§4. 超複素系とその表現類}

\paragraph{1. 拡大定理。}
いま §2 の表現定理と §3 の加群定理を組み合わせる。任意の環 \(\frR\) の代わりに、可換体 \(P\) 上の単位元を持つ超複素系 \(S,T,\ldots\) を考える：
\[
        S=x_1P+\cdots+x_nP.
\]
任意の表現環 \(\frS\) の代わりに、特に断らない限り中心 \(P\) を持つ斜体 \(A,B,\ldots\) だけを考える。この場合、\(S\) と \(A\) が元ごとに可換で、交わりが \(P\) である積環が常に存在する。これは直接積 \(S\times A\) であり、同時に \(S\) の拡大加群 \(S_A\) でもある：
\[
        S_A=x_1A+\cdots+x_nA.
\]
従ってこれを拡大環とも呼ぶ。

この特殊化のもとで、不変加群についての定理は次のようになる。

\paragraph{拡大定理。}
\(S_A\) の任意の両側イデアルは \(S\) のあるイデアルの拡大イデアル、すなわち \(S\) との交わりの拡大である。

実際、選んだ群 \(\frG\) を \(A\) の内部自己同型群とする。すると \(A\) の中心である \(P\) は完全な不変体であり、\(A\) と \(S\) の元ごとの可換性は、\(\frG\) が \(S_A\) へ拡張されて \(S\) が完全な不変領域となることを意味する。\(S_A\) の任意の両側イデアル \(\mathfrak a\) は、\(\alpha^{-1}\mathfrak a\alpha\subseteq \mathfrak a\) により許容部分群である。よってそれは交わり加群 \(\mathfrak a\cap S\) の拡大であり、この交わりは \(S\) のイデアルである。

補足：\(S\) が可換なら、\(S\) は \(S_A\) の中心である。より一般に、\(S\) の中心は \(S_A\) の中心でもある。\(A_r,B_n,\ldots\) は常に、\(A,B,\ldots\) 上の指定された次数の行列環を意味する：
\[
        A_r=\sum A c_{ik}=\sum c_{ik}A.
\]

本質的に用いるのは単純系に対する拡大定理である。

\(S\) が単純なら、\(S_A\) も単純である：
\[
        S_A=\sum Bc_{ik}=\sum c_{ik}B=B_t,
\]
ここで \(B\) は随伴する斜体であり、その中心は \(P\) を含む。ここで \(A\)、したがって \(B\) は \(P\) 上無限階数でもよい。仮定するのは \(S\) が超複素であることだけである。もし \(S\) と \(A\) が共通の中心 \(P\) を持つなら、\(P\) は \(S_A\) の中心でもある。より一般に、\(S\) が単純な超複素系なら、直接積 \(S\times A_r\) は単純である。

\paragraph{2. 表現類。}
超複素系の相反表現または直接表現とは、通常どおり、零表現ではなく、係数体 \(P\) に関して作用素準同型な表現をいう。

\paragraph{表現類に関する定理。}
\(S\) を単位元と係数体 \(P\) を持つ超複素系とし、\(A\) を中心 \(P\) を持つ斜体とする。このとき、\(S\) の \(A\) における相異なる既約相反表現類の数は、\(S_A\) をその根基で割った商環における単純右イデアル類の数に等しい。\(S\) が単純系なら、\(A\) において既約相反表現類を一つ、既約直接表現類を一つだけ持つ。

移行定理の系により、単純な相反表現類と \(S_A\) 上の単純加群類は一対一に対応する。\(S_A\) は斜体 \(A\) 上有限階数を持つので、片側イデアルについて最大条件と最小条件を満たす。\(S\) が単純なら \(S_A\) も単純である。従って単純な片側イデアル類は一つだけであり、既約相反表現類も一つだけである。相反系は直接表現についての対応する主張を与える。

\paragraph{3. 階数関係。}
\(S\) が単純で、
\[
        S_A=\sum Bc_{ik}
\]
が \(A\) 上および \(B\) 上の双方で有限階数を持つなら、
\[
        n=rt,
        \qquad n=(S:P),
        \qquad t^2=(S_A:B),                                             \tag{1}
\]
が成り立つ。ここで \(r\) は既約相反表現の次数である。

\paragraph{4. \(A\) が \(P\) 上有限階数を持つ場合の階数関係の精密化。}
この場合には三つの関係
\[
        (S:P)=n=rt,                                                     \tag{1}
\]
\[
        (B:P)\,t=(A:P)\,r,                                               \tag{2}
\]
\[
        (S:P)(B:P)=(A:P)\,r^2                                           \tag{3}
\]
がある。ここで (2) は \(S_A\) の単純右イデアルの \(P\) 上の階数を、\(B\) 上および \(A\) 上の階数で表したものである。(3) は (1) を用いて (2) に \(r\) を掛けることで従う。

\subsection*{§5. 表現定理の行列環への応用}

§4 で考えた \(S\) の \(A\) における相反表現は、\(S\) から \(A_r\) の部分環への直接準同型としても見られる。ここで \(\overline A\) は \(A\) に相反同型な斜体を表す。以下の段落の原理はその逆である。すなわち行列環 \(A_r\) の研究を、相反同型な斜体 \(A\) における表現論へ還元する。

\paragraph{1. \(A_r\) への可約および既約な埋め込み。}
\(A\) と \(\overline A\) を、その中心 \(P\) 上有限または無限階数を持つ相反同型な斜体とする。\(P\) 上の単純系 \(S\) が \(A_r\) に既約に、または可約に埋め込めるというのは、\(S\) が \(A\) において次数 \(r\) の既約、または可約な相反表現を持つときである。\(S\) が \(A_r\) に既約に埋め込めるなら、すべての行列環 \(A_{rs}\), \(s>1\), に可約に埋め込める。この同型は常に \(P\) の恒等写像を延長する。§4, 3 により、\(r\) は \(S\) の \(P\) 上の階数 \(n\) を割る。

このようにして、\(P\) を含み \(P\) 上有限階数を持つすべての単純部分環が、種々の \(A_r\) の中で得られる。実際、そのような任意の部分環は \(\overline A_r\) の部分環に相反同型であり、したがって \(A\) において可約または既約な相反表現を持つ。

\paragraph{2. 内部自己同型に関する定理。}
\(S^{(1)}\) と \(S^{(2)}\) を、\(P\) を含み \(P\) 上有限階数を持つ \(A_f\) の二つの単純部分環とする。もし \(S^{(1)}\) と \(S^{(2)}\) が \(P\) 上同型なら、この同型は \(A_f\) の内部自己同型によって誘導される。

実際、\(S^{(1)}\) と \(S^{(2)}\) は単純系 \(S\) の \(A_f\) への一般には可約な埋め込みを与え、従って \(A\) における次数 \(f\) の直接表現を与える。それらは同じ可約直接表現類に属し、変換行列が自己同型を誘導する。\(A\) 自身が \(P\) 上有限階数、従って超複素である場合には、これはよく知られた定理、すなわち中心 \(P\) を含み \(P\) 上同型な \(A_f\) の二つの単純部分系は \(A_f\) の内部自己同型によって互いに移される、という定理になる。とくに \(A_f\) の任意の自己同型は内部的である。

\paragraph{3. 元ごとに可換な部分環に関する定理。}
\(S\) を \(A_f\) に含まれ \(P\) を含む単純系とし、\(R\) を \(A_f\) の元のうち \(S\) と元ごとに可換なものすべての全体とする。このとき \(R\) も行列環である：
\[
        R=B_s.
\]
\(S_A\) の随伴斜体 \(B\) と \(R\) の随伴斜体 \(B\) は相反同型である。\(R\) と \(S\) の交わりは \(S\) の中心である。環 \(R\) が斜体であることと、\(S\) が \(A_f\) に既約に埋め込まれていることとは同値である。

§2, 3 により、環 \(R\) は、\(S\) の \(A\) における相反表現加群 \(\frM\) の自己同型環に直接同型である。これは、\(\frM\) を \(S_A\)-加群と見たときの自己同型環である。\(\frM\) が \(s\) 個の単純 \(S_A\)-加群へ分解し、§4, 1 と同じく
\[
        S_A=\sum Bc_{ik}
\]
と書けば、\(R\) は \(\sum Bc_{ik}\) に同型であり、二つの斜体は相反同型である。したがって \(R=B\) が斜体であるのは、\(s=1\)、すなわち \(S\) が既約に埋め込まれるとき、かつそのときに限る。定義により、\(R\cap S\) は \(S\) の中心である。

\paragraph{4. 超複素行列環に対する精密な交換定理。}
この場合、§4, 4 の階数関係は次の精密化を与える。

\(A\) が中心 \(P\) 上有限階数を持つとき、\(P\) を含む \(A_f\) の単純部分系は \(S,\overline S\) の対に分かれ、\(\overline S\) は \(S\) と元ごとに可換な元全体からなり、その逆も同様である。\(S_A\) と \(\overline S\) の随伴斜体は相反同型であり、\(S\) と \(\overline S_A\) の随伴斜体も同様である。\(S\) と \(\overline S\) の交わりは共通中心である。一方の部分環が斜体であることと、他方が既約に埋め込まれることとは同値である。\(S\) と \(\overline S\) の \(P\) 上の階数の積は \(A_f\) の階数である。

\(f=rs=r's'\) で、\(S\) が \(A_r\) に、\(\overline S\) が \(A_{r'}\) に既約に埋め込めるなら、\(S\) と \(\overline S\) の随伴斜体は、\(g=ss'\) として、\(A_g\) の中の可換化する部分環の一対に同型となる。

積に関する主張は階数関係から直接従う。\(S\) が既約に埋め込まれているなら、\(\overline S\) は斜体であり、\(S_A=B_t\) における \(B\) に相反同型である。§4, 4 の階数関係 (3) が主張を与える。\(S\) が可約に埋め込まれていて、\(\overline S=B_s\) なら、\(f=rs\) であり、(3) に \(s^2\) を掛けると主張を得る。これにより \(\overline S\) が \(S\) と元ごとに可換な元全体であることも示され、残りの主張は可換化する部分環に関する定理から従う。

\subsection*{§6. 単純系のガロア理論}

\paragraph{1. 中心上有限階数の斜体のガロア理論。}
\(A\) のガロア群 \(\Gcal\) は、中心 \(P\) の恒等写像を延長するすべての自己同型の群、したがって §5, 2 によりすべての内部自己同型の群として定義される。従って
\[
        \Gcal \simeq A^*/P^*,
\]
ここで \(A^*\) と \(P^*\) は \(A\) と \(P\) の非零元の乗法群を表す。したがって \(\Gcal\) の部分群 \(\mathfrak H\) は、\(P^*\) を含む \(A^*\) の部分群 \(H^*\) と
\[
        \mathfrak H\sim H^*/P^*
\]
によって一対一に対応する。\(\Gcal\) の部分群 \(\mathfrak H\) は、零を添えると \(H^*\) が斜体 \(H\) になるとき閉じていると呼ばれる。\(C\) が群 \(\mathfrak H\) を許すことは、\(C\) が \(H\) と元ごとに可換であることと同値である。従って §5, 4 の交換定理は次の形をとる。

\paragraph{非可換斜体のガロア理論の主定理。}
\(P\) と \(A\) の間の斜体 \(C\) と \(\Gcal\) の閉部分群 \(\mathfrak H\) とは、\(C\) が \(\mathfrak H\) の完全な不変斜体であり、\(\mathfrak H\) が \(C\) の完全な不変群であるように、一対一に対応する。

\paragraph{2. 単純系のガロア理論。}
\(A_f\) が中心 \(P\) を持つ単純系なら、ガロア群 \(\Gcal\) はふたたび内部自己同型群であり、従って
\[
        A_f^*/P^*
\]
に同型である。ここで \(A_f^*\) は \(A_f\) の正則元の乗法群を表す。\(\Gcal\) の部分群 \(\mathfrak H\sim H^*/P^*\) は、\(H^*\) によって生成される部分環 \(H\) が単純系で、かつ \(H^*\) が \(H\) の正則元全体であるとき、単純閉と呼ばれる。交換定理からガロア理論への移行を与えるのは次の補題である。

\paragraph{補題。}
\(P\) を含む \(A_f\) の任意の単純系は、その正則元の群 \(H^*\) によって生成される。これは \(P\) が無限個の元を持つ場合には明らかである。不定元で作った「一般元」は正則であり、適当な特殊化により \(P\) 上の正則な基底元が得られる。一般の場合にも補題は Shoda によって成り立つ。\footnote{注 3) で引用された Shoda の仕事。そこでは不定元の導入が避けられている。}

補題により、\(S\) が単純系 \(\mathfrak S\) と可換化することは、\(S\) が \(\mathfrak S\) の正則元によって誘導される自己同型を許すこととふたたび同値である。従って §5, 4 の交換定理はここでも次の形をとる。

\paragraph{単純系のガロア理論の主定理。}
\(P\) を含む \(A_f\) の単純系 \(S\) と \(\Gcal\) の単純閉部分群 \(\mathfrak H\) とは、\(S\) が \(\mathfrak H\) の完全な不変領域であり、\(\mathfrak H\) が \(S\) の完全な不変群であるように、一対一に対応する。

\paragraph{3. 拡大原理による証明。}
斜体の場合について、主定理の第二の証明を与える。それは同型、すなわち一次表現を延長する原理に基づいており、階数の数え上げを使わないので有限性仮定が少ない。とくに、無限階数の斜体 \(S\) への移行がどの方向で可能かを示す。この証明は既約相反表現類が一つだけであるという事実も用いないため、可換の場合 §8, 2 にもまったく同じように適用される。まずいくつかの補題を述べる。

\paragraph{仮定。}
\(P\) 上の単純系 \(S\) は \(A\) における一次の相反表現を許すものとし、従って斜体である。ここで \(A\) は中心 \(P\) 上有限階数でも無限階数でもよい。§4, 3 のように、\(S_A\) が \(n\) 個の単純で作用素同型な右イデアルへ分解されるものとし、
\[
        S_A=r_1+\cdots+r_n
            =e_1S_A+\cdots+e_nS_A
            =e_1A+\cdots+e_nA ;
\]
\(e_i\) によって生成される相反表現は
\[
        e_i\sigma=e_i\sigma_i
\]
により定義される。

\paragraph{補題 1.}
\(e_1,\ldots,e_n\) によって生成される \(n\) 個の相反表現は相異なる。したがって同一の表現類に属する相異なる表現である。従って \(S\) はその階数と少なくとも同数の相異なる表現を持つ。

証明のため、§2 の終わりと同じく、\(S\) からの表現を、\(A\) 上で作用素準同型な \(S_A\) の対応へ延長する。ここで対応は各 \(w\in S_A\) に対して
\[
        e_iw=e_i\omega
\]
（\(\omega\in A\)）によって定義される。もし \(e_i\) と \(e_j\), \(i\ne j\), が同じ表現を生成したなら、すべての \(w\in S_A\) について \(e_iw=e_jw\) でなければならない。これは \(e_ie_i=e_i\) かつ \(e_je_i=0\) に反する。

\paragraph{補題 2.}
\(T\) を \(P\) と \(S\) の間の斜体とし、\(s\) を \(S\) の \(T\) 上の階数とする。このとき \(T\) から \(A\) への任意の相反同型は、少なくとも \(s\) 個の異なる延長を許す。

その相反同型、すなわち \(T\) の \(A\) における相反表現が、分解
\[
        T_A=E_1T_A+\cdots+E_hT_A
            =E_1A+\cdots+E_hA,
        \qquad E_it=E_i\tau
\]
によって媒介されるとする。ここで \(E_i\) は冪等元である。これは制限ではない。基底元 \(E_i\alpha\) によって生成される表現は \(\alpha^{-1}E_i\alpha\) によって、従って冪等元によっても生成されるからである。\(S\) による右乗法は
\[
        S_A=E_1S_A+\cdots+E_hS_A
\]
を与える。各 \(E_iS_A\) は \(A\) 上階数 \(s\) を持つ。実際、\(S\) に \(T\)-基底を代入し、\(S\) と \(A\) の可換性を用いれば、\(E_iS_A\) の階数は高々 \(s\) である：
\[
        S=Tw_1+\cdots+Tw_s,
\]
\[
        E_iS_A=E_iT_Aw_1+\cdots+E_iT_Aw_s
              =E_iw_1A+\cdots+E_iw_sA .
\]
全体の和 \(S_A\) は階数 \(n=hs\) を持つので、各 \(E_iS_A\) はちょうど階数 \(s\) を持つ。従って
\[
        E_iS_A=r_{i1}+\cdots+r_{is}
              =e_{i1}A+\cdots+e_{is}A,
\]
であり、補題 1 により \(e_{i1},\ldots,e_{is}\) によって生成される \(s\) 個の表現はすべて相異なる。それらはすべて \(E_i\) によって生成される \(T\) の表現の延長である。なぜなら \(E_it=E_i\tau\) は
\[
        (e_{i1}+\cdots+e_{is})t=(e_{i1}+\cdots+e_{is})\tau
\]
を含意し、直和分解により \(e_{ij}t=e_{ij}\tau\) が従うからである。

\paragraph{定義。}
\(T\) によって誘導される、\(S\) から \(A\) への相反同型 \(\mathfrak I\) の分割 \(\mathfrak I_T\) を次のように定義する。\(\mathfrak I\) の二つの同型を、それらが \(T\) 上で同じ同型を誘導するとき、かつそのときに限って同値とみなす。

\paragraph{主定理。}
\(\mathfrak I_T\) が \(T\) によって誘導される分割なら、\(T\) はこの分割に関して極大部分斜体である。

実際、\(Z\) が \(T\) と \(S\) の間の中間斜体で、\(T\) 上の次数が \(l\) なら、補題 2 により \(\mathfrak I_T\) の各類は \(\mathfrak I_Z\) の少なくとも \(l\) 個の類に分かれる。この形から通常の形への移行は、集合 \(\mathfrak I\) の元を、固定したその一つの逆で掛けることによって得られる。集合は自己同型群となり、\(\mathfrak I_T\) の一つの類は \(T\) の不変群となる。閉部分群が完全な不変群であるという主張もこれに含まれる。\(\mathfrak H\sim H^*/P^*\) については、中間斜体 \(T\) の代わりに斜体 \(H\) を置けばよい。

\subsection*{§7. 類似代数の類。分解体}

これ以後は、中心 \(P\) 上有限階数を持つ斜体 \(A,\overline A\) と行列環 \(A_f\) だけを扱う。従ってそれらは \(P\) 上の単純正規代数であり、\(P\) 上の代数、または単に代数と呼ぶ。\(A,\overline A\) は除法代数である。同じ随伴 \(A\) を持つすべての \(A_f\) は、類似代数 \(A,A_f\) の一つの類へまとめられる。同型な \(A\) は同一視する。

\paragraph{1. 代数類の群。}
\(A,B\) が \(P\) 上の代数なら、それらの積、すなわち \(P\) 上の直接積についても同じことが成り立つ。§3 の終わりにより
\[
        A_f\times_P B_g=C_h,
\]
ここで \(C\) は中心 \(P\) を持つ除法代数で、同型を除き一意に定まる。\(A_f,B_g\) に \(P\) 上の行列環を掛けることにより、この乗法が類について一意に定まることがわかる。したがって類は乗法的に閉じた可換な系をなす。さらに R. Brauer の定理が成り立つ。類似代数の類は直接積に関してアーベル群をなす。この系は恒等類を持ち、それは \(P\) 上の行列環、すなわち \(P\) に類似な代数からなる。また各類 \((A)\) は逆類 \((\overline A)=(A)^{-1}\) を持つ。これは \(A\) に相反同型な \(\overline A\) に類似な代数からなる。このことは、§5, 3 の交換定理において \(S=A\) と特殊化することから従う。実際 \(A\) は自身に既約に埋め込めるので、\(B=P\) および \(A_A=P_f\) が得られる。\footnote{代数類の群に関するより精密な定理は因子系の理論を用いる。ここでは扱わない。ここまで、可換体に関する考察は一切現れていないことに注意せよ。例えば、階数 \((A:P)\) が平方であるという事実は用いられても証明されてもいない。}

\paragraph{2. 代数類の分解体。}
分解体の理論は、§5 の初めで述べた原理、すなわち可換拡大体を \(A\) または \(\overline A\) の中で表現するという原理に、ふたたび完全に基づく。まず次の補題を必要とする。

\paragraph{補題。}
\(A\) を中心 \(P\) を持つ斜体とし、\(Z\) を \(P\) の有限可換拡大体とする。このとき \(A_Z\) は中心 \(Z\) を持つ単純系である。実際 §4, 1 によりこれは \(A_Z\) に相反同型な系 \(Z_{\overline A}\) について成り立つ。従って、\(A\) に類似な \(P\) 上の代数類 \((A)\) に対し、体 \(Z\) は \(A_Z\) に類似な \(Z\) 上の単純代数の拡大類 \((A)_Z\) を与える。

拡大類の随伴除法代数 \(D\) は、§5, 3 の交換定理からただちに従う。

\paragraph{拡大類の随伴除法代数に関する定理。}
\((A)_Z\) を拡大類とし、\(Z\) を \(A_f\) における \(Z\) の既約埋め込み、すなわち表現とする。このとき
\[
        (A)_Z=(D),
\]
ここで \(D\) は \(A_f\) の元のうち \(Z\) と元ごとに可換なものすべての集合である。§5, 3 の単純系 \(S\) を \(Z\) に置き換え、\(Z_A\) と \(A_Z\) が相反同型であることを観察すればよい。\(Z\) が可約に埋め込まれるなら、\(Z\) と元ごとに可換な元全体は \(D\) 上の行列環である。

\(P\) の可換拡大体 \(Z\) は、拡大類 \((A)_Z\) が完全に分解する、すなわち \(Z\) 上の恒等類、つまり \(Z\) 上の全行列環の類になるとき、類 \((A)\) の分解体と呼ばれる。

\paragraph{分解体の特徴づけに関する定理。}
\(P\) の有限可換拡大体 \(Z\) が類 \((A)\) の分解体であることは、その \(A_f\) への既約埋め込みが \(A_f\) の極大可換部分体を与えることと同値である。

実際、\(Z\) が \((A)\) の分解体であるという性質は \((D)=(Z)\) と同値である。従って随伴除法代数の定理により、既約埋め込み \(Z\) は極大可換部分体でなければならない。この条件が満たされるなら、任意の \(a\in D\) を \(Z\) に添加すると可換拡大体を生じるので、必ず \(D=Z\) である。

\paragraph{注意。}
1. 同時に、極大可換部分体 \(Z\) は、既約に埋め込まれているとき、極大可換部分環を生成することが従う。なぜなら \(Z\) と元ごとに可換な元は斜体をなすからである。

2. この定理は分解体の存在も与える。随伴除法代数 \(A\) は確かに極大可換部分体を持つからである。

\paragraph{3. 階数関係、指数、無限可換拡大。}
§5, 4 の階数の主張はまず、中心上の単純代数の階数が平方であるという既知の事実を与える。\(Z\) が \(A\) の極大可換部分体なら、\(A\) へのすべての埋め込みは既約であるから、
\[
        (Z:P)(Z:P)=(A:P),
\]
従って
\[
        (A:P)=m^2,
        \qquad (Z:P)=m,
\]
ここで \(m\) は Schur の指数であり、\(A\) 自身の中に埋め込めるすべての分解体の次数でもある。一般の分解体の次数 \(n\) については
\[
        n=mr
\]
を得る。これは例えば \((A_r:P)=m^2r^2\) から、または §4, 3 の階数関係から従う。指数 \(m\) はまた、\(A_Z\) の単純成分の数でもあり、これは \(Z\) 上の階数 \(m^2\) の行列環になるからである。より一般に、\(A_Z\simeq D_l\)、\(L\) が \(A\) の \(A_r\) への既約埋め込み、\(d^2\) が中心 \(Z\) 上の除法代数 \(D\) の階数 \((D:L)\) なら、
\[
        l\,d=mr .
\]
実際、§5, 4 と 2 の随伴除法代数の定理により
\[
        (A_r:P)=(L:P)(D:P),
\]
したがって
\[
        m^2r^2=l^2d^2 .
\]

分解体の存在から、\(P\) の無限な可換、代数的または超越的拡大体 \(\Omega\) について次が従う。拡大類 \((A)_\Omega\) は中心 \(\Omega\) を持つ単純代数の類である。とくに \(\Omega\) が代数的閉体なら、\(A_\Omega\) は次数 \(m\) の行列環である。従って指数は「絶対成分数」に等しい。実際、\(Z\) が \(A\) の分解体で、\(\Omega'\) が \(\Omega\) と \(Z\) の合成体なら、\(A_{\Omega'}\) は \(\Omega'\) 上の行列環であり、従って根基を持たず、中心は \(\Omega'\) である。したがって \(A_\Omega\) も根基を持たず、その中心は \(\Omega\) である。\(A_\Omega\) の中心に属し \(\Omega\) に属さない元があれば、それは \(A_{\Omega'}\) の中心にも属してしまう。従って \(A_\Omega\) は \(\Omega\) 上の単純代数であり、\(\Omega\) が代数的閉体なら確かに \(\Omega\) 上の行列環である。\footnote{一般には「極小」分解体、すなわち真の部分体が分解体でないような分解体も、あらゆる可能な次数で存在する。注 2) で引用された Brauer--Noether を参照。}

\paragraph{4. 分離的分解体の存在。}
任意の類 \((A)\) は分離的分解体を持ち、しかも \(A\) 自身の中に埋め込めるものを持つ。\footnote{G. Köthe, Über Schiefkörper mit Unterkörpern zweiter Art über dem Zentrum, Journ. f. Math. 166 (1932), pp. 182--184 を参照。このずっと簡単な証明は M. Zorn の観察に由来する。}

基礎体 \(P\) の標数を \(p\) とし、指数を
\[
        m=s p^f
\]
と書く。ただし \(s\) は \(p\) と素である。\(Z\) が \(A\) の次数 \(m\) の分解体で、\(Z_1\) がその中に含まれる第一種拡大、すなわち
\[
        (Z:Z_1)=p^f
\]
を満たすものなら、\(A_{Z_1}\sim D\) の指数は \(p^f\) である。\(D_{Z_2}\) の指数が \(p^f\) の真の約数になるような、\(Z_1\) の分離的拡大体 \(Z_2\) の存在を示す。\(\Omega\) を代数的閉体とすると、3 により \(D_\Omega\) は根基を持たない。従って \(D\) の被約判別式は消えない。よって被約跡が非零である元 \(d\in D\) が少なくとも一つ存在する。この \(d\) はすでに基礎体 \(Z_1\) に属することはできない。なぜなら \(Z_1\) の任意の元 \(\alpha\) の跡は \(p^f\alpha\) であり、従って消えるからである。ゆえに \(Z_2=Z_1(d)\) は真の、しかも分離的な拡大体であり、\(D_{Z_2}\) の指数は \(p^f\) の真の約数である。これを有限回繰り返すと、次数 \(m\) が \(A\) の指数に一致する \((A)\) の分離的分解体が得られる。

\subsection*{§8. 可換の場合における分解体、部分分解体、ガロア理論}

中心上単純な系の分解体から任意の単純系の分解体へ移るためには、なお中心の分解理論を展開し、それに §6, 3 と平行なガロア理論を加えなければならない。この段落に現れるすべての体と系は可換である。

\paragraph{1. 可換単純系の分解体と部分分解体。}
可換系 \(Z\) を \(P\) 上単純、従って体とし、\(\Omega\) を \(P\) の代数的閉拡大体とする。\(Z_\Omega\) が根基を持たない系であり続けることは、\(Z\) が \(P\) 上分離的、すなわち第一種拡大であることと同値であることが知られている。このときに限り、\(P\) 上の階数と同じだけの一次同型写像を \(\Omega\) へ持ち、従って同じだけの相異なる表現を持つ。この場合それらは表現類と一致する。

\(Z_\Omega\) に根基が現れる可能性、従って絶対単純成分への直和分解が必ずしも起こらない可能性は、次の定義へ導く。\(P\) の拡大体 \(\Lambda\) は、\(Z_\Lambda\) が一次の組成因子へ分解する、言い換えればすべての絶対既約表現がすでに \(\Lambda\) にあるとき、分解体と呼ばれる。\(P\) の拡大体 \(T\) は、\(Z_T\) が少なくとも一つの一次組成因子を分離する、言い換えれば \(Z\) の少なくとも一つの絶対既約表現が \(T\) に存在するとき、部分分解体と呼ばれる。

分解体と部分分解体は次の定理によって特徴づけられる。

\paragraph{定理。}
体 \(T\) が部分分解体であることは、\(T\) が \(Z\) に同型な部分体 \(Z'\) を含むことと同値である。体 \(\Lambda\) が分解体であることは、\(Z\) に属するガロア体に同型な部分体 \(\Gamma\) を含むことと同値である。

これは、絶対既約表現による分解体と部分分解体の定義から直接従う。とくに非可換の場合との類比で次の系を得る。体 \(Z'\) が極小部分分解体、すなわち真の部分体が部分分解体でないものであることは、\(Z'\) が \(Z\) に同型であることと同値である。極小部分分解体が同時に分解体であることは、\(Z\) が正規、従って \(P\) 上ガロア体であることと同値である。

これが、単純代数をその中心上正規と呼ぶ理由である。\(P\) 上の固定された代数的閉体 \(\Omega\) の中では、極小分解体はただ一つ、すなわち \(Z\) に属するガロア体だけである。これは、非可換の場合に一般に無限個の非同型体があることと対照的である。その理由は、可換の場合の直和分解の一意性にあり、非可換の場合にはそれが作用素同型を除いてしか一意でないことにある。同じことを別の形でいえば、可換の場合には相異なる絶対既約表現が有限個しかないのに対し、非可換の場合には同じ類に属していても相異なるものが無限個あるからである。

\paragraph{2. 分離的体 \(Z/P\) の場合のガロア理論の超複素的証明。}
§6, 3 と正確に同じ類比で、任意の \(P\) 上分離的な \(Z\) について同型の理論が成り立つ。\(Z\) がガロアなら、通常の自己同型群へ移ることができる。

\paragraph{仮定。}
\(P\) 上の単純系 \(Z\) は有限分離拡大体であり、\(\Omega\) は \(P\) 上有限または無限の分解体、\(Z'=Z^{(1)}\) は \(Z\) に同型な部分体、\(\Gamma\) は対応するガロア体を表す。1 により直和分解
\[
        Z_\Omega=r^{(1)}+\cdots+r^{(n)}
        =e^{(1)}Z_\Omega+\cdots+e^{(n)}Z_\Omega
        =e^{(1)}\Omega+\cdots+e^{(n)}\Omega
\]
がある。\(e^{(i)}\) によって生成される表現 \(Z\to Z^{(i)}\), ただし \(Z^{(i)}\subseteq\Gamma\), は
\[
        e^{(i)}z=e^{(i)}z^{(i)}
\]
により定義される。

\paragraph{補題 1.}
\(e^{(1)},\ldots,e^{(n)}\) によって生成される \(n\) 個の表現はすべて相異なる。証明は §6, 3 と同じであり、また 1 からも従う。それらは相異なる表現類に属する。

\paragraph{補題 2.}
\(T\) が \(P\) と \(Z\) の間の中間体で、\(s\) が \(Z\) の \(T\) 上の階数なら、\(T\) から \(\Omega\) への任意の同型は少なくとも、従って正確に、\(s\) 個の延長を許す。

証明は §6, 3 と同じである。「少なくとも」が「正確に」へ強められるのは、\(Z_\Omega\) の分解の一意性により、\(Z\) が \(\Omega\) の中に少なくともではなく正確に \(n\) 個の同型を持つからである。

§6, 3 と同じく、\(Z\) の同型からなる有限集合 \(\mathfrak I\) の、\(T\) によって誘導される分割 \(\mathfrak I_T\) を定義する。すなわち、\(\mathfrak I\) のうち \(T\) 上で同じ同型を誘導するものを、ちょうど同値とみなす。すると次を証明できる。

\paragraph{主定理。}
\(\mathfrak I_T\) が \(T\) によって誘導される分割なら、\(T\) はこの分割に関する極大部分体である。

ここから、\(Z\) がガロアである場合には、§6, 3 とまったく同じように、同型の逆と合成することによって自己同型群と通常の定式化へ移る。部分群についての対応する主張は、冪等元の考察から従う。この考察はそれ自体にも意義がある。

\paragraph{3. \(Z_\Omega\) の冪等元。}
\(S\) を、\(Z\) を個々の \(Z^{(i)}\) へ送る \(n\) 個の置換とする。つまり、\(Z\to Z_1\) と \(Z\to Z^{(i)}\) の同型の合成であり、前者は \(Z\to Z_1\) の相反である。体 \(Z\) はガロアである必要はない。\(a\in Z_\Omega\) に対し、置換 \(S\) は \(Z\) 上恒等的であると約束することにより作用素として定義される。したがって各 \(a\) について、\(Z_{(i)}\) に属する共役 \(a^S\) が定義される。この約束のもとで次の定理が成り立つ。

\paragraph{定理。}
\(Z_\Omega\) の \(n\) 個の冪等元 \(e^{(i)}\) は共役である：
\[
        e^{(i)}=e^S,
        \qquad e=e^{(1)};
\]
それらは \(n\) 個の共役な拡大系 \(Z\Omega\) に属する。

実際、\(S\) を作用させると、定義関係
\[
        ez=ez^{(1)},
        \qquad e^2=e
\]
は
\[
        e^S z=e^S z^{(i)},
        \qquad (e^S)^2=e^S
\]
へ移る。分解の一意性により冪等元自身も一意に定まるので、\(e^S=e^{(i)}\) である。さらに \(Z\) は部分分解体でもあり、表現 \(ez=ez^{(1)}\) はすでに媒介されているから、\(e\) はすでに \(Z\Omega\) に属し、従って \(e^S\) は \(Z^S\Omega\) に属する。

\(Z\) がガロアなら、ガロア理論の主定理のうち部分群に関する部分はただちに従う。\(\mathfrak H\) がガロア群 \(\Gcal\) の部分群なら、直和の成分の一意性により、元
\[
        \sum_{H\in\mathfrak H} e^H
\]
は \(\mathfrak H\) のすべての置換を許し、それ以外は許さない。群 \(\Gcal\) は \(Z\) に対して定義されたので、これは \(Z\) のガロア理論であり、同型により \(Z'\) の理論も決定する。

\paragraph{4. 冪等元と補基底との関係。}
ここでも \(Z\) がガロアであるとは仮定しない。

\paragraph{定理。}
\(a_1,\ldots,a_n\) を \(P\) 上の \(Z\) の基底とし、冪等元を
\[
        e=a_1\beta_1+\cdots+a_n\beta_n
\]
と書く。すると
\[
        e^S=a_1\beta_1^S+\cdots+a_n\beta_n^S .
\]
このとき
\[
        \beta_1^S,\ldots,\beta_n^S
\]
は、\(e^S\) によって媒介される写像のもとで、\(a_1,\ldots,a_n\) に相補的な基底 \(b_1,\ldots,b_n\) の像である。

\(e^S\) によって媒介される写像では \(e^Sx=e^Sx^S\)、とくに \(e^Sa_i=e^Sa_i^S\) であり、関係
\[
        e^Se^S=e^S,
        \qquad e^Se^R=0\quad(S\ne R)
\]
は割当て \(e^S\mapsto1\), \(e^R\mapsto0\) を与える。従って
\[
        1=a_1^S\beta_1^S+\cdots+a_n^S\beta_n^S,
        \qquad
        0=a_1^R\beta_1^S+\cdots+a_n^R\beta_n^S\quad(S\ne R).
\]
これを行列で書くと、相補基底の定義式
\[
\begin{pmatrix}
 a_1^1 & \cdots & a_n^1\\
 \vdots & \ddots & \vdots\\
 a_1^n & \cdots & a_n^n
\end{pmatrix}
\begin{pmatrix}
 \beta_1^S\\ \vdots\\ \beta_n^S
\end{pmatrix}
=
\begin{pmatrix}
0\\ \vdots\\ 1\\ \vdots\\ 0
\end{pmatrix}
\]
である。跡による定義への移行は、
\[
        a_i=\sum_S a_i^S e^S,
        \qquad b_i=\sum_S b_i^S e^S
\]
について、行を入れ替えた後の行列関係が
\[
        \delta_{ij}=\Sp(a_i b_j)
\]
になることから従う。ここで \(a_i\) と \(b_i\) はすべての置換 \(S\) を許すので \(Z\) に属する。\footnote{表現論 §25 を参照。}\footnote{相補基底と冪等元との関係は Dedekind, Zur Theorie der aus \(n\) Haupteinheiten gebildeten komplexen Größen に初めて現れる。全集第 II 巻、pp. 4--5 を参照。}

相補基底の反変性は、ここでは冪等元 \(e^S\) が \(Z\) の不変元であり、選ばれた基底に依存しないことから従う。したがって、
\[
        a_1\beta_1^S+\cdots+a_n\beta_n^S
\]
はすべて、\(Z/P\) の別の基底
\[
        a_1',\ldots,a_n'
\]
へ移っても自分自身に変換される。

\subsection*{§9. 任意の系の分解体と部分分解体}

\paragraph{1. 定義と単純な場合への還元。}
§8 で与えた定義に、一般の場合には次の定義が対応する。\(S\) を \(P\) 上超複素とする。\(P\) の可換拡大体 \(\Lambda\) は、片側イデアルによる組成列において \(S_\Lambda\) が絶対単純な組成因子へ分解するとき、\(S\) の分解体と呼ばれる。言い換えれば、\(S\) の \(\Lambda\) におけるすべての既約表現がすでに絶対既約であるときである。\(P\) の拡大体 \(T\) は、\(S_T\) が少なくとも一つの絶対単純組成因子を分離するとき、同値にいえば、\(T\) において既約な少なくとも一つの表現がすでに絶対既約であるとき、部分分解体と呼ばれる。

これらの定義は、可換の場合について §8 で、単純代数について §7 で与えられた定義を明らかに含む。後者については、中心の任意の可換拡大に対して \(A_\Lambda\) が完全可約であり、従って組成因子と直和分解の成分が一致するからである。中心上の全行列環、そしてそれだけが、絶対単純成分を与え、従って絶対既約表現も与える。さらに \(A_\Lambda\) は両側単純であり、作用素同型な成分へ分解する。従って一つの絶対単純成分を分離することはすでに完全分解を与え、部分分解体と分解体は一致する。

定義からただちに次の事実が従う。\(S\) の分解体は、\(P\) において既約な各表現に対して一つずつ分解体を取った合成体である。\(S\) の部分分解体は、\(P\) において既約な表現の任意の部分分解体である。単純系、すなわち根基で割った剰余類環の成分と、\(P\) において既約な表現との間には一対一対応があるので、単純系を考えれば十分である。その場合の結果は §7 と §8 を組み合わせて得られる。

\paragraph{2. 単純系の分解体と部分分解体。}
まず、\(P\) 上分離的な中心 \(Z\) を持つ単純系 \(A\) の場合を考える。§8, 1 および §8, 3 から次が得られる。\(\Omega\) を \(Z\) の分解体とすると、\(Z_\Omega\) の直和分解に対応する \(A_\Omega\) の両側分解は、\(A\) にそれぞれ同型な \(n\) 個の共役成分
\[
        e^S A_\Omega
\]
を与え、それらはそれぞれ中心
\[
        e^S Z_\Omega=e^S Z
\]
上の単純代数となる。置換 \(S\) は、\(A\) 上恒等的であると約束することにより、\(A_\Omega\) に対する作用素として定義される。

実際、§8, 3 により冪等元は共役であるから、\(A\) に対する置換の定義により成分 \(e^S A_\Omega\) も共役である。\(A\) は両側単純なので、\(e^S A_\Omega\) は環として \(A\) に同型である。さらに \(e^SA_\Omega=e^S A_{Z^S}\) である。なぜなら \(e^S Z_\Omega=e^S Z\) だからである。従って中心は係数領域と一致し、成分は正規代数である。

§7, 2 の結果からさらに次が従う。\(A\) が \(P\) 上の単純系で、中心 \(Z\) が \(P\) 上分離的なら、任意の有限または無限拡大体 \(\Omega\) に対して \(A_\Omega\) も根基を持たず、従って完全可約である。実際、§7, 2 により \(e^S A_\Omega\) は任意の係数拡大の後も単純であり、根基を持たない。従って直和についても同じである。この直和は \(A_\Omega\) であり、また \(\Omega\) が \(Z\) の分解体でないときには \(A_\Omega\) から係数拡大によって得られる。

§7, 2 から次の特徴づけも従う。

\paragraph{部分分解体と分解体の特徴づけ。}
\(A\) が分離的な中心 \(Z\) を持つ除法代数なら、部分分解体の既約埋め込みは、\(Z\) を含む \(A_f\) の極大可換部分体のすべて、かつそれだけによって与えられる。分解体は、部分分解体と中心の分解体、特に中心に属するガロア体との合成体のすべて、かつそれだけである。極小部分分解体は、中心がガロアでなくても分解体であり得る。

部分分解体は \(Z\) に同型な体を含まなければならない。部分分解体の埋め込みに関する主張は、前の事実と §7, 2 から従う。さらに分解体は \(Z\) の分解体 \(\Omega\) を含まなければならない。しかし \(\Omega\) 上の任意の部分分解体は同時に分解体である。なぜなら成分 \(e^S A_\Omega\) は単純であり、\(\Omega\) に同型な中心を持つからである。\(Z\) がガロアなら、極小部分分解体と分解体は一致する。しかし非ガロアの場合にも、可換の場合とは違って、同時に分解体である極小部分分解体が存在し得る。すなわち、そのような極小部分分解体が \(Z\) に属するガロア体を含む場合である。例：中心が \(P(\sqrt2)\) に同型な四元数除法代数。ただし \(P\) は有理数体とする。このとき \(P(\sqrt2)\) に属するガロア体は極小部分分解体かつ分解体である。

最後に、中心が非分離的な場合には、\(Z_\Omega\) も、従って \(A_\Omega\) も根基を持つ系となる。\(\mathfrak C\) を \(Z_\Omega\) の根基、\(\mathfrak C A_\Omega\) を \(A_\Omega\) における拡大イデアルとする。すると \(Z_\Omega/\mathfrak C\) は、§8, 2 と正確に同じ類比で、共役な成分への直和分解を持つ。これにより上と同じく \(A_\Omega/\mathfrak C A_\Omega\) の直和分解が得られ、その成分 \(e^SA_\Omega\) は \(A\) に同型で、中心は \(e^SZ^\Omega\) である。従って部分分解体と分解体に関する定理はなお成り立つ。さらに階数の数え上げから、\(\mathfrak C A_\Omega\) の組成列に対応する組成因子は、単なる準同型ではなく、\(A_\Omega/\mathfrak C A_\Omega\) と作用素同型であることがわかる。

既約表現について述べると、次のまとめになる。\(P\) 上既約な超複素系の表現が与えられているとき、その表現が絶対的に完全可約であることは、対応する中心が \(P\) 上分離的であることと同値である。いずれの場合にも、その表現の分解体と部分分解体は、対応する単純系への埋め込みによって上のように特徴づけられる。とくに、部分分解体の最低次数は、中心の次数と中心上の対応する代数類の指数との積であり、任意の部分分解体の次数はその倍数である。表現は、組成列の意味で、中心に含まれる最大分離体の次数だけの、相異なるが共役な絶対既約表現類へ分解する。各類は、指数と、この分離拡大上の中心の次数との積で与えられる回数だけ現れる。

完全基礎体の場合、すなわち分離拡大体だけが現れる場合には、これは I. Schur の既知の結果を返す。ただし、R. Brauer にすでに現れる埋め込みによる特徴づけが付け加わっている。

\begin{center}
(1932年6月8日受理。)
\end{center}
